\section{Clases laterales, normalidad y teoremas de isomorfismo}\label{sec:grupos-cocientes}

\subsection{Subgrupos normales y cocientes de grupos}\label{tema:5-1-1}\label{tema:5-1-2}
\begin{defi}
Un subgrupo $N\le G$ es \emph{normal}, escrito $N\trianglelefteq G$, si satisface
\[
gng^{-1}\in N\qquad(g\in G,\ n\in N).
\]
\end{defi}
Si $G$ es abeliano, cada subgrupo es normal, pues $gng^{-1}=n$. Para un subconjunto $A\subseteq G$ usamos la notación
\[
gA=\set{ga\given a\in A},\qquad
Ag=\set{ag\given a\in A}.
\]
La normalidad admite las siguientes formulaciones equivalentes:
\begin{align*}
N\trianglelefteq G
&\quad\Longleftrightarrow\quad gNg^{-1}\subseteq N\quad\text{para todo }g\in G\\
&\quad\Longleftrightarrow\quad gNg^{-1}=N\quad\text{para todo }g\in G\\
&\quad\Longleftrightarrow\quad gN=Ng\quad\text{para todo }g\in G.
\end{align*}
En efecto, aplicar la inclusión a $g^{-1}$ da $g^{-1}Ng\subseteq N$ y, al conjugar por $g$, obtenemos $N\subseteq gNg^{-1}$. Multiplicar a la derecha por $g$ transforma la igualdad de conjugados en la igualdad de clases laterales. Incluso la condición $gN\subseteq Ng$ para todo $g$ basta: equivale a $gNg^{-1}\subseteq N$.

\begin{lema}\label{lema:C5-nucleo-normal}
El núcleo de todo homomorfismo $\varphi:G\to G'$ es un subgrupo normal de $G$.
\end{lema}
\begin{proof}
Ya sabemos que es un subgrupo. Si $n\in\ker\varphi$ y $g\in G$, entonces
\begin{align*}
\varphi(gng^{-1})
&=\varphi(g)\varphi(n)\varphi(g^{-1})\\
&=\varphi(g)e_{G'}\varphi(g)^{-1}=e_{G'}.
\end{align*}
Así $gng^{-1}\in\ker\varphi$.
\end{proof}

Si $\sim$ es una relación de equivalencia en el conjunto $G$, la proyección $\pi:G\to G/{\sim}$ envía $g$ a $[g]$. Toda función $f:G\to X$ constante en clases induce una única función $\bar f([g])=f(g)$:
\[
\begin{tikzcd}[column sep=large,row sep=large]
G/{\sim}\arrow[r,destacado,"\exists!\,\bar f"]&X\\
G\arrow[u,two heads,"\pi"]\arrow[ur,"f"']&
\end{tikzcd}
\]
La función está bien definida porque $g\sim g'$ implica $f(g)=f(g')$. La sobreyectividad de $\pi$ da la unicidad. Es decir, $(G/{\sim},\pi)$ es inicial entre las funciones desde $G$ que identifican elementos equivalentes, con flechas dadas por las factorizaciones correspondientes.

Queremos ahora una operación de grupo en el cociente para la cual $\pi$ sea un homomorfismo. Necesariamente debe ser
\[
[a][b]=[ab].
\]
El punto decisivo es que el resultado no dependa de los representantes.

\begin{prop}\label{prop:C5-compatible}
La regla $[a][b]=[ab]$ define una estructura de grupo en $G/{\sim}$ si y sólo si, para todos $a,a',g\in G$,
\begin{equation}\label{eq:C5-compatible}
a\sim a'\quad\Longrightarrow\quad
ga\sim ga'\quad\text{y}\quad ag\sim a'g.
\end{equation}
En ese caso, $\pi$ es un homomorfismo y todo homomorfismo $\varphi:G\to G'$ constante en clases induce un único homomorfismo $\bar\varphi:G/{\sim}\to G'$ con $\bar\varphi\circ\pi=\varphi$.
\end{prop}
\begin{proof}
Si la regla define una operación, $[a]=[a']$ implica
\[
[ga]=[g][a]=[g][a']=[ga'],\qquad
[ag]=[a][g]=[a'][g]=[a'g].
\]
Recíprocamente, supongamos \eqref{eq:C5-compatible}. Si $a\sim a'$ y $b\sim b'$, entonces $ab\sim a'b\sim a'b'$, así que $[ab]=[a'b']$. La operación está bien definida. La asociatividad se hereda de $G$, el neutro es $[e_G]$ y el inverso de $[g]$ es $[g^{-1}]$. La igualdad $\pi(ab)=\pi(a)\pi(b)$ es la propia regla de multiplicación.

La propiedad universal de conjuntos produce $\bar\varphi([a])=\varphi(a)$. Es un homomorfismo, pues
\[
\bar\varphi([a][b])=\varphi(ab)
=\varphi(a)\varphi(b)=\bar\varphi([a])\bar\varphi([b]).
\]
La unicidad como función implica la unicidad como homomorfismo.
\end{proof}
Decimos que una equivalencia que satisface \eqref{eq:C5-compatible} es \emph{compatible con la estructura de grupo}.
\begin{obs}
La igualdad de clases no exige igualdad de sus representantes. Por ejemplo, en $\mathbb Z/2\mathbb Z$, $[0]=[2]$ aunque $0\ne2$.
\end{obs}

\subsection{Equivalencias invariantes y clases laterales}\label{tema:5-2-1}
\begin{prop}\label{prop:C5-clases}
Sea $\sim$ una equivalencia en $G$ invariante por multiplicación a la izquierda. Entonces $H=[e_G]$ es un subgrupo y
\[
a\sim b\quad\Longleftrightarrow\quad a^{-1}b\in H
\quad\Longleftrightarrow\quad aH=bH.
\]
En particular, $[a]=aH$.
\end{prop}
\begin{proof}
El conjunto $H$ contiene a $e_G$. Si $a,b\in H$, de $a\sim e_G$ y $b\sim e_G$ obtenemos $a\sim b$. Multiplicando a la izquierda por $a^{-1}$, resulta $e_G\sim a^{-1}b$, por lo que $a^{-1}b\in H$. Esta variante del criterio de subgrupo también basta: tomando $a=b$ se obtiene el neutro; tomando $b=e_G$ se obtienen inversos; y, una vez disponibles, $(a^{-1})^{-1}b=ab\in H$.

Para $a,b$ cualesquiera, multiplicar $a\sim b$ por $a^{-1}$ da $e_G\sim a^{-1}b$, y multiplicar esta última relación por $a$ da el recíproco. Si $a^{-1}b\in H$, escribimos $b=ah$ con $h\in H$; entonces $bH=ahH=aH$, porque $hH=H$. Finalmente, si $aH=bH$, de $b\in aH$ se deduce $b=ah$ para algún $h\in H$, luego $a^{-1}b\in H$. La descripción de la clase de $a$ se sigue de la misma equivalencia.
\end{proof}

\begin{defi}
Para $H\le G$, los conjuntos $aH$ se llaman \emph{clases laterales izquierdas}, y los conjuntos $Ha$, \emph{clases laterales derechas}.
\end{defi}

\begin{prop}\label{prop:C5-correspondencias}
Los subgrupos de $G$ se corresponden biyectivamente con las equivalencias invariantes por multiplicación a la izquierda, mediante
\[
H\longmapsto\bigl(a\sim_L b\ \Longleftrightarrow\ a^{-1}b\in H\bigr).
\]
También se corresponden con las equivalencias invariantes por multiplicación a la derecha, mediante
\[
H\longmapsto\bigl(a\sim_R b\ \Longleftrightarrow\ ab^{-1}\in H\bigr).
\]
Las clases de estas relaciones son $aH$ y $Ha$, respectivamente.
\end{prop}
\begin{proof}
Para $\sim_L$, reflexividad, simetría y transitividad se siguen de
\[
a^{-1}a=e\in H,\qquad
b^{-1}a=(a^{-1}b)^{-1},\qquad
a^{-1}c=(a^{-1}b)(b^{-1}c).
\]
Además, $(ga)^{-1}(gb)=a^{-1}b$, lo que da la invariancia por la izquierda. La clase del neutro es $H$ y las demás son $aH$. La proposición anterior muestra que toda equivalencia de este tipo se recupera de su clase del neutro; ambas construcciones son inversas.

Para $\sim_R$, usamos $aa^{-1}=e$, $ba^{-1}=(ab^{-1})^{-1}$ y $ac^{-1}=(ab^{-1})(bc^{-1})$, junto con $(ag)(bg)^{-1}=ab^{-1}$. Sus clases son $Ha$: por ejemplo, $ab^{-1}\in H$ equivale a $a=hb$ para algún $h\in H$, y por tanto a $Ha=Hb$. El mismo argumento, con izquierda y derecha intercambiadas, prueba que la clase del neutro de una equivalencia invariante por la derecha es un subgrupo y recupera la relación.
\end{proof}
Para todo $h\in H$ tenemos $hH=Hh=H$. Siempre $a\in aH\cap Ha$; si una clase izquierda $aH$ coincide con una clase derecha $Hb$, entonces, por contener a $a$, esta última es $Ha$, y $aH=Ha$.

\begin{ejem}\label{ejem:C5-S3}
En $S_3$, sean $H=\{e,(1\ 2)\}$ y $g=(1\ 3\ 2)$. Con nuestra convención de composición, el factor derecho actúa primero y
\[
gH=\{(1\ 3\ 2),(2\ 3)\},\qquad
Hg=\{(1\ 3\ 2),(1\ 3)\}.
\]
Son diferentes: $g(1\ 2)$ es la transposición $(2\ 3)$, mientras que $(1\ 2)g$ es la transposición $(1\ 3)$.
\end{ejem}

\subsection{Normalidad y construcción del cociente}\label{tema:5-2-2}
\begin{prop}\label{prop:C5-normal-relaciones}
Las relaciones $\sim_L$ y $\sim_R$ asociadas a $H\le G$ coinciden si y sólo si $H$ es normal.
\end{prop}
\begin{proof}
Las relaciones coinciden exactamente cuando tienen las mismas clases para cada elemento, es decir, cuando $gH=Hg$ para todo $g\in G$. Esta condición equivale a $H\trianglelefteq G$.
\end{proof}
Si $H$ es normal, su relación de clases laterales es invariante tanto por la izquierda como por la derecha. La \cref{prop:C5-compatible} permite entonces definir lo siguiente.

\begin{defi}
Para $H\trianglelefteq G$, el \emph{grupo cociente} es el conjunto $G/H=\set{gH\given g\in G}$ con la operación
\[
(aH)(bH)=(ab)H.
\]
Su neutro es $H=e_GH$ y el inverso de $gH$ es $g^{-1}H$.
\end{defi}
La proyección $\pi:G\to G/H$, $g\mapsto gH$, es un homomorfismo sobreyectivo, y
\[
\ker\pi=\set{g\in G\given gH=H}=H.
\]
Así, los subgrupos normales son exactamente los núcleos de homomorfismos de grupos.

\begin{teo}[Propiedad universal del cociente]\label{teo:C5-universal}
Si $H\trianglelefteq G$ y $\varphi:G\to G'$ es un homomorfismo con $H\subseteq\ker\varphi$, existe un único homomorfismo $\bar\varphi:G/H\to G'$ tal que $\bar\varphi\circ\pi=\varphi$:
\[
\begin{tikzcd}[column sep=large,row sep=large]
G\arrow[r,"\varphi"]\arrow[d,two heads,"\pi"']&G'\\
G/H\arrow[ur,destacado,"\exists!\,\bar\varphi"']&
\end{tikzcd}
\]
\end{teo}
\begin{proof}
Si $aH=bH$, entonces $a^{-1}b\in H\subseteq\ker\varphi$ y
\[
\varphi(a)^{-1}\varphi(b)=\varphi(a^{-1}b)=e_{G'}.
\]
Por ello $\varphi(a)=\varphi(b)$. El homomorfismo es constante en clases, y la \cref{prop:C5-compatible} produce el único mediador, dado por $\bar\varphi(gH)=\varphi(g)$.
\end{proof}
Por normalidad de $H$, las condiciones $ab^{-1}\in H$ y $a^{-1}b\in H$ son equivalentes.

\begin{ejem}
Como $\mathbb Z$ es abeliano, $n\mathbb Z$ es normal. En notación aditiva,
\[
a\sim b\quad\Longleftrightarrow\quad b-a\in n\mathbb Z
\quad\Longleftrightarrow\quad n\mid(b-a).
\]
Las clases son $[a]_n=a+n\mathbb Z$, y la operación del cociente satisface $[a]_n+[b]_n=[a+b]_n$. Recuperamos así el grupo de clases de congruencia introducido en la sección~\ref{tema:2-2-2}.
\end{ejem}

\subsection{Índice y recuento de clases laterales}\label{tema:5-2-3}
El \emph{índice} $[G:H]$ es el número de clases laterales izquierdas de $H$. Para hablar del conjunto $G/H$ no hace falta que $H$ sea normal; la normalidad es necesaria para su estructura de grupo cociente.

\begin{prop}[Lagrange]\label{prop:C5-lagrange}
Si $G$ es finito y $H\le G$, entonces $|G|=[G:H]|H|$. En particular, $|H|$ divide a $|G|$.
\end{prop}
\begin{proof}
Las clases laterales forman una partición de $G$, y cada una tiene $|H|$ elementos: la función $H\to gH$, $h\mapsto gh$, es biyectiva, con inversa $x\mapsto g^{-1}x$. Sumamos sus cardinalidades sobre las $[G:H]$ clases.
\end{proof}
Si $H\le K\le G$ y $[G:H]$ es finito, también
\begin{equation}\label{eq:C5-indices}
[G:H]=[G:K][K:H].
\end{equation}
Para verlo, elijamos representantes $g_i$ de las clases de $K$ en $G$ y $k_j$ de las de $H$ en $K$. Los $g_ik_j$ representan todas las clases de $H$ en $G$. Si $g_ik_jH=g_{i'}k_{j'}H$, primero $g_iK=g_{i'}K$, luego $i=i'$; al cancelar $g_i$, se sigue $k_jH=k_{j'}H$ y $j=j'$. Así hay exactamente el producto indicado de clases, incluso si los grupos son infinitos.

\subsection{Primer teorema de isomorfismo y productos}\label{tema:5-3-0}
\begin{teo}[Primer teorema de isomorfismo]\label{teo:C5-primer-iso}
Si $\varphi:G\to G'$ es un homomorfismo, la siguiente aplicación inducida es un isomorfismo:
\[
\widetilde\varphi:G/\ker\varphi\xrightarrow{\ \sim\ }\im\varphi,
\qquad g\ker\varphi\longmapsto\varphi(g).
\]
En consecuencia, se factoriza como
\[
\begin{tikzcd}[column sep=large]
G\arrow[r,two heads,"\pi"]\arrow[rrr,bend left=22,"\varphi"]&
G/\ker\varphi\arrow[r,destacado,"\widetilde\varphi","\sim"']&
\im\varphi\arrow[r,hook,"i"]&G'.
\end{tikzcd}
\]
\end{teo}
\begin{proof}
El núcleo es normal, de modo que existe el cociente. La propiedad universal, aplicada al homomorfismo con codominio restringido a $\im\varphi$, produce $\widetilde\varphi$. Es sobreyectivo por definición de la imagen. Si dos clases tienen la misma imagen, entonces $\varphi(a)=\varphi(b)$, de donde
\[
\varphi(a^{-1}b)=e_{G'},\qquad a^{-1}b\in\ker\varphi.
\]
Así $a\ker\varphi=b\ker\varphi$ y $\widetilde\varphi$ es inyectivo. Es un homomorfismo biyectivo y, por la \cref{prop:C3-iso}, es un isomorfismo. Las fórmulas muestran $i\circ\widetilde\varphi\circ\pi=\varphi$.
\end{proof}
\begin{coro}\label{coro:C5-iso-sobreyectivo}
Si $\varphi:G\to G'$ es sobreyectivo, entonces $G/\ker\varphi\cong G'$.
\end{coro}
\begin{proof}
En ese caso $\im\varphi=G'$ y aplicamos el teorema anterior.
\end{proof}

\begin{ejem}[Cociente de un producto]
Si $H_1\trianglelefteq G_1$ y $H_2\trianglelefteq G_2$, entonces
\[
H_1\times H_2\trianglelefteq G_1\times G_2,\qquad
\frac{G_1\times G_2}{H_1\times H_2}
\cong\frac{G_1}{H_1}\times\frac{G_2}{H_2}.
\]
La aplicación se obtiene usando las proyecciones $p_i$ y los cocientes $q_i$:
\[
\begin{tikzcd}[column sep=large,row sep=large]
&G_1\arrow[r,two heads,"q_1"]&G_1/H_1\\
G_1\times G_2\arrow[ur,"p_1"]\arrow[dr,"p_2"']\arrow[rr,destacado,"\Pi"]&&
(G_1/H_1)\times(G_2/H_2)\arrow[u]\arrow[d]\\
&G_2\arrow[r,two heads,"q_2"']&G_2/H_2
\end{tikzcd}
\]
Por la propiedad universal del producto, $\Pi$ es el único homomorfismo cuyas componentes son $q_1p_1$ y $q_2p_2$. En concreto,
\[
\Pi(g_1,g_2)=(g_1H_1,g_2H_2).
\]
Es sobreyectivo, pues cada par de clases tiene un par de representantes. Además,
\begin{align*}
\ker\Pi
&=\set{(g_1,g_2)\given g_1H_1=H_1,\ g_2H_2=H_2}\\
&=\set{(g_1,g_2)\given g_1\in H_1,\ g_2\in H_2}
=H_1\times H_2.
\end{align*}
Como núcleo, este producto es normal; el \cref{coro:C5-iso-sobreyectivo} da el isomorfismo anunciado. Su fórmula explícita es
\[
(g_1,g_2)(H_1\times H_2)\longmapsto(g_1H_1,g_2H_2).
\]
\end{ejem}

Tomando $H_1=\{e_{G_1}\}$ y $H_2=G_2$, resulta
\[
\frac{G_1\times G_2}{\{e_{G_1}\}\times G_2}
\cong (G_1/\{e_{G_1}\})\times(G_2/G_2)
\cong G_1\times\{*\}\cong G_1.
\]
El isomorfismo es la proyección sobre la primera coordenada, pasada al cociente. La abreviatura $(G_1\times G_2)/G_2$ utiliza la inclusión $g_2\mapsto(e_{G_1},g_2)$.

En particular, $C_6\cong C_2\times C_3$ y, al cocientar por el subgrupo de orden $3$, obtenemos $C_6/C_3\cong C_2$. Para hacer explícita esta identificación, el homomorfismo
\[
\mathbb Z/6\mathbb Z\to(\mathbb Z/2\mathbb Z)\times(\mathbb Z/3\mathbb Z),
\qquad [a]_6\mapsto([a]_2,[a]_3),
\]
es inyectivo porque su núcleo exige que $2$ y $3$ dividan a $a$, y por tanto que $6$ lo divida. Ambos grupos tienen seis elementos, así que es un isomorfismo. El subgrupo $\ang{[2]_6}$ corresponde a $\{[0]_2\}\times\mathbb Z/3\mathbb Z$, y el cociente conserva la primera coordenada.

\subsection{Correspondencia y teoremas de isomorfismo}\label{tema:8-1-1b}
\begin{lema}[Correspondencia e isomorfismos de cocientes]\label{lema:C8-cocientes}
Sean $N\trianglelefteq G$ y $\pi:G\to G/N$ la proyección.
\begin{enumerate}
\item Los subgrupos de $G/N$ corresponden biyectivamente a los subgrupos de $G$ que contienen a $N$, por $\bar H\mapsto\pi^{-1}(\bar H)$ y $H\mapsto H/N$. La correspondencia preserva inclusiones y normalidad en los respectivos grupos ambientes.
\item Si $A\le G$ y $B\trianglelefteq G$, entonces $AB\le G$ y
\[
A/(A\cap B)\cong AB/B.
\]
Si además $A\trianglelefteq G$, entonces $AB\trianglelefteq G$.
\item Si $N\subseteq L\trianglelefteq K\le G$, entonces
\[
(K/N)/(L/N)\cong K/L.
\]
\end{enumerate}
\end{lema}
\begin{proof}
La preimagen de un subgrupo es subgrupo. Si $N\subseteq H$, entonces $\pi^{-1}(\pi(H))=H$: la igualdad $\pi(g)=\pi(h)$ implica $g\in hN\subseteq H$. La sobreyectividad de $\pi$ da $\pi(\pi^{-1}(\bar H))=\bar H$. Conjugar clases equivale a proyectar conjugados; por ello $H\trianglelefteq G$ equivale a $H/N\trianglelefteq G/N$. El mismo razonamiento vale al restringir la proyección a cualquier subgrupo que contenga a $N$.

Para el segundo punto, $ba=a(a^{-1}ba)\in AB$ y $a^{-1}ba\in B$, así que los productos e inversos de elementos de $AB$ vuelven a estar allí. La aplicación $A\to AB/B$, $a\mapsto aB$, es un homomorfismo sobreyectivo con núcleo $A\cap B$. El primer teorema de isomorfismo da el resultado. Si $A$ es también normal, $g(AB)g^{-1}=(gAg^{-1})(gBg^{-1})=AB$.

Finalmente, $kN\mapsto kL$ define un homomorfismo sobreyectivo $K/N\to K/L$: está bien definido porque $N\subseteq L$. Su núcleo es $L/N$, y se aplica otra vez el primer teorema de isomorfismo.
\end{proof}
